\let\negmedspace\undefined
\let\negthickspace\undefined
\documentclass[journal]{IEEEtran}
\usepackage[a5paper, margin=10mm, onecolumn]{geometry}
%\usepackage{lmodern} % Ensure lmodern is loaded for pdflatex
\usepackage{tfrupee} % Include tfrupee package

\setlength{\headheight}{1cm} % Set the height of the header box
\setlength{\headsep}{0mm}     % Set the distance between the header box and the top of the text

\usepackage{gvv-book}
\usepackage{gvv}
\usepackage{cite}
\usepackage{amsmath,amssymb,amsfonts,amsthm}
\usepackage{algorithmic}
\usepackage{graphicx}
\usepackage{textcomp}
\usepackage{xcolor}
\usepackage{txfonts}
\usepackage{listings}
\usepackage{enumitem}
\usepackage{mathtools}
\usepackage{gensymb}
\usepackage{comment}
\usepackage[breaklinks=true]{hyperref}
\usepackage{tkz-euclide} 
\usepackage{listings}
% \usepackage{gvv}                                        
\def\inputGnumericTable{}                                 
\usepackage[latin1]{inputenc}                                
\usepackage{color}                                            
\usepackage{array}                                            
\usepackage{longtable}                                       
\usepackage{calc}                                             
\usepackage{multirow}                                         
\usepackage{hhline}                                           
\usepackage{ifthen}                                           
\usepackage{lscape}
\begin{document}
\bibliographystyle{IEEEtran}
\title{12.536}
\author{EE25BTECH11002 - Achat Parth Kalpesh }
{\let\newpage\relax\maketitle}
\renewcommand{\thefigure}{\theenumi}
\renewcommand{\thetable}{\theenumi}
\setlength{\intextsep}{10pt} % Space between text and floats
\numberwithin{equation}{enumi}
\numberwithin{figure}{enumi}
\renewcommand{\thetable}{\theenumi}
\parindent 0px



\textbf{Question:}\\
The rank of the matrix 
$
\myvec{
1 & 1 & 0 & -2 \\
2 & 0 & 2 & 2 \\
4 & 1 & 3 & 1 \\
}$
is

\begin{multicols}{4}
\begin{enumerate}
    \item 1
    \item 2
    \item 3
    \item 4
\end{enumerate}
\end{multicols}

\textbf{Solution:}\\
We perform row operations to reduce it to row echelon form:


\begin{align}
\myvec{
1 & 1 & 0 & -2 \\
2 & 0 & 2 & 2 \\
4 & 1 & 3 & 1
}
&\xleftrightarrow[]{R_2 \leftarrow R_2 - 2R_1}
\myvec{
1 & 1 & 0 & -2 \\
0 & -2 & 2 & 6 \\
4 & 1 & 3 & 1
} \\
\myvec{
1 & 1 & 0 & -2 \\
0 & -2 & 2 & 6 \\
4 & 1 & 3 & 1
}&\xleftrightarrow[]{R_3 \leftarrow R_3 - 4R_1}
\myvec{
1 & 1 & 0 & -2 \\
0 & -2 & 2 & 6 \\
0 & -3 & 3 & 9
} \\
\myvec{
1 & 1 & 0 & -2 \\
0 & -2 & 2 & 6 \\
0 & -3 & 3 & 9
}&\xleftrightarrow[]{R_2 \leftarrow -\frac{1}{2} R_2}
\myvec{
1 & 1 & 0 & -2 \\
0 & 1 & -1 & -3 \\
0 & -3 & 3 & 9
} \\
\myvec{
1 & 1 & 0 & -2 \\
0 & 1 & -1 & -3 \\
0 & -3 & 3 & 9
}&\xleftrightarrow[]{R_1 \leftarrow R_1 - R_2}
\myvec{
1 & 0 & 1 & 1 \\
0 & 1 & -1 & -3 \\
0 & -3 & 3 & 9
} \\
\myvec{
1 & 0 & 1 & 1 \\
0 & 1 & -1 & -3 \\
0 & -3 & 3 & 9
}&\xleftrightarrow[]{R_3 \leftarrow R_3 + 3 R_2}
\myvec{
1 & 0 & 1 & 1 \\
0 & 1 & -1 & -3 \\
0 & 0 & 0 & 0
}
\end{align}

Now the matrix is in reduced row echelon form \\
The number of non-zero rows = 2

\begin{align}
    \text{rank}\brak{A} = 2
\end{align}


\end{document}